% ============================================================
% D04 — Hak Akses dan Menu
% ============================================================
\flowmeta{D04}{Hak Akses dan Menu}{Umum}

\begin{flowsection}{Tujuan}
Menjelaskan peran pengguna dan bagaimana menu Panel Admin ditampilkan sesuai
peran (\emph{role-based menu access}).
\end{flowsection}

\begin{flowsection}{Peran Pengguna}
Sistem mengenal beberapa peran berikut:
\begin{center}
\renewcommand{\arraystretch}{1.4}
\begin{tabular}{|l|l|l|}
\hline
\textbf{Peran} & \textbf{Akses} & \textbf{Keterangan} \\ \hline
\texttt{admin} & Panel Admin & Akses penuh ke modul ERP sesuai menu peran \\ \hline
\texttt{warehouse\_staff} & Panel Admin & Operasional gudang/stok \\ \hline
\texttt{finance\_officer} & Panel Admin & Keuangan, kas, dan laporan \\ \hline
\texttt{customer} & Storefront & Berbelanja sebagai pelanggan \\ \hline
\end{tabular}
\end{center}
\end{flowsection}

\begin{flowsection}{Menu Dinamis}
Daftar menu Panel Admin tidak ditulis tetap di kode, melainkan diambil dari
basis data (\emph{master menu} + pemetaan peran-ke-menu). Saat admin login,
sistem mengirim daftar menu yang boleh diakses oleh perannya. Bila sebuah peran
tidak memiliki menu, sisi kiri tidak menampilkan menu apa pun dan URL admin
tersebut tidak dapat diakses.
\end{flowsection}

\shot{gambar/umum/d04-1.png}{Sidebar menu Panel Admin setelah login (menu tampil sesuai peran admin)}{fig:d04-1}

\begin{flowsection}{Langkah Melihat Hak Akses}
\begin{enumerate}
  \item Login sebagai admin pada halaman \texttt{/login}.
  \item Setelah berhasil, buka \texttt{/admin/dashboard}.
  \item Amati sidebar: daftar menu muncul sesuai peran akun tersebut.
  \item Menu \textbf{Master Pengguna}, \textbf{Peran}, dan \textbf{Kelola Menu}
        (bagian Pengaturan) mengatur peran dan hak akses menu.
\end{enumerate}
\end{flowsection}

\begin{flowsection}{Hasil yang Diharapkan}
Admin melihat menu sesuai perannya; halaman yang tidak diizinkan tidak muncul di
menu dan tidak dapat diakses langsung lewat URL.
\end{flowsection}

\begin{flowsection}{Penanganan Error}
Jika muncul pesan ``Tidak ada menu yang tersedia untuk peran Anda'', berarti akun
tersebut belum memiliki pemetaan menu. Admin perlu menetapkan menu untuk peran
tersebut melalui modul \textbf{Role--Menu Mapping}.
\end{flowsection}

\begin{flowsection}{Kaitan Modul}
FE \texttt{AdminSidebar.jsx}, \texttt{App.jsx} (guard view admin). BE
\texttt{AuthController@menus}, \texttt{MenuController}, rute \texttt{/api/auth/menus}.
Tabel \texttt{roles}, \texttt{menus}, \texttt{role\_menus}.
\end{flowsection}

\docver{v1.0}{Sep 2026}
